% SOURCE_AUTHORITY: sga5_fr_workpass.tex, lines 5431--5478 (LF-normalized unit).
% SOURCE_SHA256: 791F4EFFC5E02832D5D77ED03518C8156D6F07E4C8238B03545DB93D883FBB28
\subsubsection*{5.12. Traços e produtos tensoriais externos}

Sejam, para $i=1,2$, $A_i$ uma $K$-álgebra plana e $A=A_1\otimes_KA_2$. Sejam $E_i\in\ob D^-(A_i)$, $F_i\in\ob D^b(A_i^e)$ tais que $\uRHom_{A_i}(E_i,F_i)\in\ob D^b(A_i^\circ)$, $F=F_1\lten_KF_2\in\ob D^b(A^e)$, $\uRHom_A(E,F)\in\ob D^b(A^\circ)$, onde $E=E_1\lten_KE_2$. Tem-se então um quadrado comutativo
\begin{equation}\tag{5.12.1}
\begin{tikzcd}[column sep=large,row sep=large]
(\uRHom_{A_1}(E_1,F_1)\lten_{A_1}E_1)\lten_K(\uRHom_{A_2}(E_2,F_2)\lten_{A_2}E_2) \arrow[r,"(1)"] \arrow[d,"(2)"'] & (F_1\lten_{A_1^e}A_1)\lten_K(F_2\lten_{A_2^e}A_2) \arrow[d,"(3)"]\\
\uRHom_A(E,F)\lten_AE \arrow[r] & F\lten_{A^e}A,
\end{tikzcd}
\end{equation}
onde (1) é o produto tensorial das flechas 5.3.5, (2) está definida pela flecha de produto tensorial
\begin{equation}\tag{5.12.2}
\uRHom_{A_1}(E_1,F_1)\lten_K\uRHom_{A_2}(E_2,F_2)\longrightarrow\uRHom_A(E,F),
\end{equation}
e (3) é o isomorfismo canônico evidente. A verificação é essencialmente trivial. Do mesmo modo se verifica que, para $E_i\in\ob D^-(A_i)$, $F_i\in\ob D^+(A_i^e)$, $G_i\in\ob D^-(A_i)$, com $F_i\lten_{A_i}G_i\in\ob D^b(A_i)$, $\uRHom_{A_i}(E_i,F_i)\in\ob D^b(A_i^\circ)$, $\uRHom_{A_i}(E_i,F_i\lten_{A_i}G_i)\in\ob D^b(K)$, $F\lten_AG\in\ob D^b(A)$ (onde $G=G_1\lten_KG_2$), tem-se um quadrado comutativo
\begin{equation}\tag{5.12.3}
{\scriptstyle
\begin{tikzcd}[column sep=2.8em,row sep=large]
\begin{gathered}
(\uRHom_{A_1}(E_1,F_1)\lten_{A_1}G_1)\lten_K{}\\[-.15em]
(\uRHom_{A_2}(E_2,F_2)\lten_{A_2}G_2)
\end{gathered}
\arrow[r,"(1)"] \arrow[d,"(2)"']
&
\begin{gathered}
\uRHom_{A_1}(E_1,F_1\lten_{A_1}G_1)\lten_K{}\\[-.15em]
\uRHom_{A_2}(E_2,F_2\lten_{A_2}G_2)
\end{gathered}
\arrow[d,"(3)"]\\
\uRHom_A(E,F)\lten_AG \arrow[r,"5.4.2"'] & \uRHom_A(E,F\lten_AG),
\end{tikzcd}
}
\end{equation}
onde (2) e (3) são flechas de produto tensorial do tipo 5.12.2, e (1) é o produto tensorial das flechas 5.4.2. Da comutatividade dos quadrados 5.12.1 e 5.12.3 resulta que, para $E_i\in\ob D(A_i)_{\mathrm{parf}}$, $F_i\in\ob D^b(A_i^e)$ tais que $F=F_1\lten_KF_2\in\ob D^b(A^e)$, tem-se um quadrado comutativo
\begin{equation}\tag{5.12.4}
\begin{tikzcd}[column sep=large,row sep=large]
\uRHom_{A_1}(E_1,F_1\lten_{A_1}E_1)\lten_K\uRHom_{A_2}(E_2,F_2\lten_{A_2}E_2) \arrow[r,"(1)"] \arrow[d,"(2)"'] & (F_1\lten_{A_1^e}A_1)\lten_K(F_2\lten_{A_2^e}A_2) \arrow[d,"(3)"]\\
\uRHom_A(E,F\lten_AE) \arrow[r,"5.6.1"'] & F\lten_{A^e}A,
\end{tikzcd}
\end{equation}
onde (1) é o produto tensorial das flechas $\Tr_{A_i}$ (5.6.1), (2) está definida por 5.12.2 e (3) é o isomorfismo canônico que figura em 5.12.1. Por conseguinte, para $u_i\in\uHom_{A_i}(E_i,F_i\lten_{A_i}E_i)$, tem-se
\begin{equation}\tag{5.12.5}
\Tr_A(u_1\lten_Ku_2)=\Tr_{A_1}(u_1)\Tr_{A_2}(u_2),
\end{equation}
onde o produto do segundo membro é definido pela flecha canônica
\[
H^0(T,F_1\lten_{A_1^e}A_1)\otimes H^0(T,F_2\lten_{A_2^e}A_2)
\longrightarrow H^0(T,F\lten_{A^e}A).
\]
